30--50\,Hz vibration applied to the neck or ear is proposed here as a simple vagal patch: a localized vibrotactile interface intended to engage afferent vagal tone through mechanically delivered stimulation. In the RBM framework, the design premise is not that vibration ``treats'' a condition by itself, but that a narrow frequency band may provide a practical entry point for modulating autonomic state through resonance-sensitive pathways.

A minimal implementation would use a wearable or adhesive transducer positioned at the cervical region or near the ear, where stimulation can be delivered comfortably and repeatedly. The target band, 30--50\,Hz, is selected as a working range for vagus-related vibration protocols and should be treated as an experimental parameter rather than a universal prescription. In practice, the device would emphasize stable contact, low amplitude, and controlled duty cycle so that the stimulus remains localized and tolerable.

Within the broader prototype set, the vagal patch serves as the most direct example of a low-barrier, body-mounted frequency intervention: a compact hardware form factor, a clearly defined frequency window, and a plausible autonomic target. Its role in the RBM program is to provide a testable bridge between mechanical vibration and measurable changes in inflammation-related or parasympathetic markers, especially heart-rate variability and subjective calm.
